% =============================================================================
\section{Sinais e sistemas contínuos: vocabulário e grandezas}
% =============================================================================

\subsection*{Teoria mínima}

Um \textbf{sinal} é uma função do tempo; no domínio contínuo escrevemos $x(t)$,
$t\in\mathbb{R}$. Três famílias resolvem praticamente todos os problemas de PDS:

\begin{itemize}
  \item \textbf{Exponencial complexa:} $x(t)=e^{st}$, $s=\sigma+j\Omega$. Se
        $\sigma=0$ temos a senoide $e^{j\Omega t}=\cos\Omega t+j\operatorname{sen}\Omega t$.
  \item \textbf{Degrau:} $u(t)=1$ para $t\ge 0$ e $0$ para $t<0$; o ``liga/desliga''
        que torna os sinais causais.
  \item \textbf{Impulso:} $\delta(t)$, definido por $\int x(t)\delta(t-t_0)\,dt=x(t_0)$.
\end{itemize}

\paragraph{Energia e potência média.}
\[
  E=\int_{-\infty}^{\infty}|x(t)|^2\,dt, \qquad
  P=\lim_{T\to\infty}\frac{1}{T}\int_{-T/2}^{T/2}|x(t)|^2\,dt .
\]
Sinal de energia: $0<E<\infty$ e $P=0$. Sinal de potência: $0<P<\infty$ e
$E=\infty$. Para a senoide real $x(t)=A\cos(\Omega_0 t+\phi)$ vale sempre
$P=A^2/2$, \emph{independentemente} de $\Omega_0$ e $\phi$.

\paragraph{Periodicidade.} $x(t)$ é periódico se existe $T>0$ com
$x(t+T)=x(t)$. Para uma soma $x(t)=\sum_i A_i\cos(\Omega_i t+\phi_i)$, a
periodicidade ocorre se e somente se todas as razões $\Omega_i/\Omega_j$ são
\emph{racionais}; o período fundamental é $T_0=1/f_0$, com $f_0=\mathrm{mdc}$
das frequências em Hz.

\paragraph{Classificação de sistemas.} Um sistema $y=\mathcal{T}\{x\}$ é
\begin{itemize}
  \item \textbf{sem memória} se $y(t)$ depende só de $x(t)$;
  \item \textbf{causal} se $y(t)$ depende apenas de $x(\tau)$, $\tau\le t$;
  \item \textbf{invariante no tempo (TI)} se deslocar a entrada desloca a saída;
  \item \textbf{linear} se $\mathcal{T}\{a x_1+b x_2\}=a\,\mathcal{T}\{x_1\}+b\,\mathcal{T}\{x_2\}$;
  \item \textbf{estável (BIBO)} se entrada limitada $\Rightarrow$ saída limitada.
\end{itemize}
Linear $+$ invariante no tempo $=$ \textbf{SLIT}, a classe que domina a
disciplina e cuja análise se resume à convolução (Capítulo~\ref{sec:slit}).

\paragraph{Formas de onda de referência.}
\[
  \operatorname{rect}(t)=
  \begin{cases}1,&|t|<1/2\\0,&|t|>1/2\end{cases},
  \qquad
  \operatorname{tri}(t)=
  \begin{cases}1-|t|,&|t|<1\\0,&|t|>1\end{cases},
  \qquad
  \operatorname{sinc}(t)=\frac{\operatorname{sen}(\pi t)}{\pi t}.
\]

\subsection*{Exercícios resolvidos}

% ---------------------------------------------------------------- 1.1
\exer{$x(t)=3\cos(200\pi t+\pi/4)$. Determine amplitude, frequência
(rad/s e Hz), período, valor de pico e potência média.}

\textbf{Solução.} Comparando com $A\cos(\Omega_0 t+\phi)$:
\[
  A=3,\qquad \Omega_0=200\pi\ \text{rad/s},\qquad
  f_0=\frac{\Omega_0}{2\pi}=100\ \text{Hz},\qquad
  T_0=\frac{1}{f_0}=0{,}01\ \text{s}=10\ \text{ms}.
\]
Valor de pico $=A=3$ (unidade do sinal). Potência média:
\[
  P=\frac{A^2}{2}=\frac{9}{2}=4{,}5\ \text{W}.
\]

\resp{$A=3$, $\Omega_0=200\pi$ rad/s, $f_0=100$ Hz, $T_0=10$ ms, $P=4{,}5$ W.}

% ---------------------------------------------------------------- 1.2
\exer{Calcule a energia de $x(t)=e^{-2t}u(t)$ e classifique o sinal.}

\textbf{Solução.} Aplicando a definição e integrando:
\[
  E=\int_0^{\infty}e^{-4t}\,dt
   =\left[-\frac{e^{-4t}}{4}\right]_{0}^{\infty}
   =\frac{1}{4}=0{,}25\ \text{J}.
\]
Como $0<E<\infty$, o sinal é de \textbf{energia} (logo $P=0$). Generalizando,
$x(t)=e^{-at}u(t)$ tem $E=1/(2a)$ para $a>0$ — resultado usado várias vezes
adiante.

\resp{$E=1/(2a)=0{,}25$ J; sinal de energia.}

% ---------------------------------------------------------------- 1.3
\exer{$x(t)=2\cos(4\pi t)+3\operatorname{sen}(6\pi t)$ é periódico? Em caso
afirmativo, determine o período fundamental e a potência média.}

\textbf{Solução.} As frequências são $f_1=2$ Hz e $f_2=3$ Hz; a razão
$f_1/f_2=2/3$ é \textbf{racional}, portanto a soma é periódica. O período
fundamental é o menor $T$ com $2T\in\mathbb{Z}$ \emph{e} $3T\in\mathbb{Z}$:
\[
  f_0=\mathrm{mdc}(2,3)=1\ \text{Hz}\ \Rightarrow\ T_0=\frac{1}{f_0}=1\ \text{s}.
\]
Potência: como $\cos$ e $\operatorname{sen}$ de frequências distintas são
ortogonais em um período, as potências somam:
\[
  P=\frac{2^2}{2}+\frac{3^2}{2}=2+4{,}5=6{,}5\ \text{W}.
\]

\resp{$T_0=1$ s, $P=6{,}5$ W.}

\begin{obs}
Se a razão das frequências fosse irracional (por exemplo $f_2=1/(2\pi)$ Hz),
a soma \emph{não} seria periódica, mesmo sendo limitada.
\end{obs}
